% (User input window)

\paragraph{边界作为压力函数的非解析点：谱半径、维数谱与相变的统一坐标}
定理 \ref{thm:pom-real-q-pressure-analytic-interpolation} 将整数 \(q\) 的 moment 指纹提升为一条实参数压力曲线，并把压力固化为一个谱对象：
\[
\tau(q)=\log\lambda(q)=\log\rho(\mathcal L_q).
\]
在该框架内，“边界”可被制度化为压力的非解析前沿；其机制并非来自几何图像的直观突变，而来自主谱点的可竞争性（谱隙坍缩或主谱分支切换）对解析扰动制度的破坏。

\begin{definition}[压力边界集合（非解析前沿）]\label{def:pom-pressure-nonanalytic-boundary}
沿用定理 \ref{thm:pom-real-q-pressure-analytic-interpolation} 的记号，在开区间 \(J\subset(0,\infty)\) 上令
\(
\tau(q):=\log\rho(\mathcal L_q)
\)
为实参数压力。
定义压力边界集合为
\[
\partial\tau:=\{q\in J:\ \tau\ \text{在}\ q\ \text{处不为实解析}\}.
\]
\end{definition}

\begin{remark}[谱隙制度下的解析性与边界触发机制]\label{rem:pom-nonanalyticity-as-spectral-gap-closure}
若 \(q_0\in J\) 满足定理 \ref{thm:pom-real-q-pressure-analytic-interpolation} 的一致谱隙假设（在某邻域内 Perron 主特征值为代数重数 \(1\) 的孤立谱点），则由谱投影的解析扰动理论，\(\lambda(q)\) 与 \(\tau(q)\) 在 \(q_0\) 的邻域内实解析。
因此 \(\partial\tau\) 的出现只能由谱结构退化触发：典型机制包括谱隙塌缩、主谱点失去孤立性、或出现等模竞争导致主谱分支切换。
\end{remark}

\begin{remark}[维数谱的谱半径化：与“分形维数—矩阵—相变”的对齐]\label{rem:pom-dimension-matrix-pressure-unified}
维数侧的两个基本接口在文内均已给出。
其一，在前缀超度量下子移位的 Hausdorff 维数由拓扑熵决定，从而对有限图/邻接算子退化为谱半径的对数刻度（推论 \ref{cor:Ym-amb-hausdorff-codim}）。
其二，对折叠推前测度 \(w_m=d_m/2^m\) 的 Rényi 维数谱，在几何标度 \(\varepsilon_m=\varphi^{-m}\) 下由
\[
D_q=\frac{q\log 2-\log r_q}{(q-1)\log\varphi}
\]
给出闭式（命题 \ref{prop:pom-renyi-dimension-spectrum}），并且在整数锚点上满足 \(\tau(q)=\log r_q\)（定理 \ref{thm:pom-real-q-pressure-analytic-interpolation}）。
因此，“维数谱的结构转折”可在同一坐标系中等价转写为 \(\tau\) 的谱机制变化：当 \(q\) 方向上出现压力非解析或谱隙显著变弱时，由 \(\tau\) 诱导的多重分形 Legendre 几何与大偏差谱等号（定理 \ref{thm:pom-ldp-uniform-and-wm-exact-spectrum}）将同步出现可审计的边界效应。
在附录 \ref{app:pom-projective-pressure-operator} 的射影转移算子口径下，\(\tau\) 的可微性/曲率与谱隙可在 \(q\approx 4\) 的邻域被直接审计；若出现尖点、曲率峰或谱隙塌缩，则给出“边界＝相变”的谱证书。
\end{remark}

